\section{政策建议}

\subsection{高龄照护缺口区：补护理型服务}

\subsubsection{区域特征}

高龄人口密度高，医疗和护理服务不足。

\subsubsection{补位方向}

新增护理型社区养老服务点、长护险定点服务、社区康复和上门护理。

\subsubsection{实施建议}

优先依托社区卫生服务中心和现有养老服务站升级。


\subsection{助餐服务不足区：补长者食堂和社区食堂}

\subsubsection{区域特征}

老年需求高，但长者食堂和助餐点不足。

\subsubsection{补位方向}

新增长者食堂、社区食堂、移动助餐和配送服务。

\subsubsection{实施建议}

优先布局在老旧小区和高需求网格周边。


\subsection{居家照料不足区：补日间照料中心和养老服务站}

\subsubsection{区域特征}

机构养老不一定缺，但社区居家服务弱。

\subsubsection{补位方向}

新增日间照料中心、社区养老服务站、助浴助洁助医服务。

\subsubsection{实施建议}

推动服务站向综合照护站升级。


\subsection{代际共享服务短板区：补复合型公共服务空间}

\subsubsection{区域特征}

养老服务需求高，同时缺少公园、体育空间、图书馆、社区活动中心等共享服务。

\subsubsection{补位方向}

补充社区食堂、公共活动空间、文体设施和便民服务中心。

\subsubsection{实施建议}

推动“一站多用”的复合型社区服务空间建设。


\subsection{供给冗余区：推动存量设施转型}

\subsubsection{区域特征}

服务供给较多，但老年需求相对不足。

\subsubsection{优化方向}

控制新增设施，提升服务质量，推动存量设施向助餐、康复、文体和上门服务转型。

\subsubsection{实施建议}

避免重复建设，提高现有设施利用率。
